\documentclass[11pt]{article}
    \usepackage[utf8]{inputenc}
    \usepackage[T1]{fontenc}
    \usepackage[margin=1in]{geometry}
    \usepackage{hyperref}
    \usepackage{enumitem}
    \title{Indi: Giving directions & transition word}
    \author{Piter Garcia}
    \date{March 20, 2026}
    \begin{document}
    \maketitle
    \section*{Quick Scan}
    \begin{itemize}[leftmargin=*]
    \item Grade: Grade 2
    \item Workbook sessions: Session 29
    \item Class blocks: Day 1 10:45-11:35 Callon, Day 1 2:15-3:05 Henchen, Day 4 10:45-11:35 Barthelman
    \item Reference file: indi-giving-directions-and-transition-word.bib
    \end{itemize}
    \section*{School Planning Goals and Inclusion}
    \begin{itemize}[leftmargin=*]
    \item What is expected: I can give clear directions, using transitions words to get from one place to another. | | I can use details (from my partners directions) to perform a task. | | I can develop and document a plan to get Indi to a certain location.
    \item What we achieved: This lesson points to local transcript or evidence files in the workspace so the packet can be revised from what actually happened in class.
    \item What needs improvement: stronger proof, cleaner assessment notes, and tighter lesson artifacts.
    \item How to improve: add transcript proof, one saved artifact, and clearer success criteria.
    \item Inclusion focus: use visual modeling, chunked directions, predictable routines, and flexible participation roles.
    \end{itemize}
    \section*{Official References}
    \begin{itemize}[leftmargin=*]
    \item \url{https://csteachers.org/k12standards/interactive/}
\item \url{https://dmmedia.sphero.com/email-marketing/Sphero-Edu/SpheroEdu-k12-teacher-resource-guide-v1_updated050818.pdf}

    \end{itemize}
    \end{document}
